%! AP Statistics — Unit 3, Lesson 3.3 — Homework
\documentclass[10pt]{article}
\usepackage{apstats-article}
\usepackage{apstats-boxes}
\usepackage{grading}
\gradingassignment{apstats-u03-l03-hw}
\gradingcourse{advanced_statistics}

\begin{document}

\pageheader{Unit 3, Lesson 3.3}{Homework}
{\small\itshape Write each final answer in its box. Only what is inside a box is graded.}\par

\vspace{0.1in}

\begin{practicebox}
\small
\textbf{1.} Identify the sampling method in each scenario.

\vspace{6pt}
\begin{enumerate}[leftmargin=1.5em, itemsep=10pt]
  \item A university has 20{,}000 students. A researcher assigns each student a
        number and uses a random number generator to select 400.\\
        Method: \makebox[7cm][l]{\vphantom{\blank{7cm}}\hspace{46bp}\smash{\gradedblank[id=hw1a, type=short, points=1, answer={simple random sample}, check=text]{2.4in}{0.2in}}}

  \item A county divides its population into 10 geographic zones. Researchers
        randomly select 3 zones and survey every household in those zones.\\
        Method: \makebox[7cm][l]{\vphantom{\blank{7cm}}\hspace{46bp}\smash{\gradedblank[id=hw1b, type=short, points=1, answer={cluster sample}, check=text]{2.4in}{0.2in}}}

  \item A polling agency calls every 50th person on a voter registration list,
        starting with a randomly selected number between 1 and 50.\\
        Method: \makebox[7cm][l]{\vphantom{\blank{7cm}}\hspace{46bp}\smash{\gradedblank[id=hw1c, type=short, points=1, answer={systematic random sample}, check=text]{2.4in}{0.2in}}}

  \item A car manufacturer divides its workers into three job categories (assembly,
        quality control, shipping), then randomly selects 20 workers from each.\\
        Method: \makebox[7cm][l]{\vphantom{\blank{7cm}}\hspace{46bp}\smash{\gradedblank[id=hw1d, type=short, points=1, answer={stratified random sample}, check=text]{2.4in}{0.2in}}}
\end{enumerate}
\end{practicebox}

\vspace{0.1in}

\begin{practicebox}
\small
\textbf{2.} Stratified vs.\ Cluster

A school district has 5 elementary schools (each with about 400 students).
A researcher wants to sample 100 students to ask about their reading habits.

\begin{enumerate}[leftmargin=1.5em, itemsep=8pt]
  \item Describe how to use a \emph{stratified} sample (strata = schools):
        \gradedhide{\writeline\\[1pt]\gradedhere{\gradedblank[id=hw2a, type=rubric, points=2, display={One stratum per school: an SRS of 20 students from each of the 5 schools.}]{\linewidth+6pt}{2in}}\writeline\\[1pt]\writelines{9}}
  \item Describe how to use a \emph{cluster} sample (clusters = schools):
        \gradedhide{\writeline\\[1pt]\gradedhere{\gradedblank[id=hw2b, type=rubric, points=2, display={Randomly select 1 school (or enough for about 100 students); survey every student in it.}]{\linewidth+6pt}{2in}}\writeline\\[1pt]\writelines{9}}
  \item Are the schools better described as strata or clusters? Explain.
        \gradedhide{\writeline\\[1pt]\gradedhere{\gradedblank[id=hw2c, type=rubric, points=2, display={Clusters: each school is a mix of all kinds of students (heterogeneous), not a homogeneous group.}]{\linewidth+6pt}{2in}}\writeline\\[1pt]\writelines{9}}
\end{enumerate}
\end{practicebox}

\vspace{0.1in}

\begin{practicebox}
\small
\textbf{3.} Free Response:\\
A state health agency wants to study the physical activity levels of high school
students. Describe a \emph{stratified random sample} you could use. Identify the
strata, explain why stratifying is appropriate, and describe the selection procedure.

\vspace{6pt}
\gradedhide{\gradedhere{\gradedblank[id=hw3, type=rubric, points=3, display={Strata: grade levels 9--12, since activity likely differs by grade (e.g., required PE in 9th). From each grade's enrollment list, randomly select students in proportion to its size.}]{\linewidth+6pt}{2in}}\writeline\\[1pt]\writelines{9}}
\end{practicebox}

\end{document}
